\title{DeepSeekMath: Pushing the Limits of Mathematical Reasoning in Open Language Models}

\begin{abstract}
The task of mathematical reasoning remains a formidable obstacle for language models, given its highly structured and intricate demands. This study presents DeepSeekMath~7B, an advancement based on further pre-training of DeepSeek-Coder-Base-v1.5 7B with a dataset comprising 120B math-centric tokens derived from Common Crawl, as well as textual and programming data. Remarkably, DeepSeekMath~7B attains a score of 51.7\% on the MATH benchmark, which features competition-level problems, all without dependence on auxiliary toolkits or ensemble methods, thereby nearing the capabilities exhibited by Gemini-Ultra and GPT-4. When evaluating with 64 self-consistency samples, DeepSeekMath~7B reaches 60.9\% accuracy on MATH. The advanced mathematical reasoning skills of DeepSeekMath~can be traced to two primary innovations: firstly, the effective utilization of web-scale public data through a rigorously designed data selection process; and secondly, the development of Group Relative Policy Optimization (GRPO)—a modified version of Proximal Policy Optimization (PPO)—that boosts reasoning proficiency while simultaneously streamlining PPO's memory consumption.
\end{abstract}

\section{Introduction}

Large language models (LLM) have dramatically transformed methods for mathematical reasoning in AI, propelling notable progress in both quantitative reasoning benchmarks \citep{MATH} and geometry-focused evaluation datasets \citep{trinh2024solving}. Beyond this, such models serve as valuable aids to humans faced with intricate mathematical challenges \citep{tao}. Nonetheless, state-of-the-art systems like GPT-4 \citep{gpt4} and Gemini-Ultra \citep{gemini} remain unavailable to the public, and open-source alternatives continue to substantially lag these closed models in terms of mathematical performance.

This paper presents DeepSeekMath, a language model tailored to mathematical domains that not only surpasses existing open-source models in mathematical reasoning, but also attains near-GPT-4-level performance on academic evaluation tasks. Central to this advancement is the construction of the DeepSeekMath Corpus, a meticulously curated pre-training resource comprising 120B mathematics-specific tokens. The data is sourced from Common Crawl (CC), with relevant content identified through a classifier based on fastText \citep{joulin2016fasttext}. Initially, the classifier is trained with positive samples from OpenWebMath \citep{openwebmath} and a diverse assortment of unrelated web content as negative samples. This classifier is subsequently used to extract further positive samples from CC, after which these are rigorously screened through human review. The refined dataset then enables retraining of the classifier, enhancing its accuracy. Empirical results suggest that this expansive dataset is of considerable quality; for example, the foundational model DeepSeekMath-Base 7B attains scores of 64.2\% on GSM8K \citep{gsm8k} and 36.2\% on the advanced MATH dataset \citep{MATH}, surpassing the performance of Minerva 540B \citep{minerva}. The multilingual nature of the DeepSeekMath Corpus also yields measurable gains on Chinese mathematics benchmarks \citep{wei2023cmath,agieval}. We propose that the methodologies developed for mathematical data curation here offer a baseline for future research, while noting that substantial further enhancements are possible.

DeepSeekMath-Base is built upon DeepSeek-Coder-Base-v1.5 7B \citep{deepseek-coder}, as empirical evidence suggests that initializing from a code-focused pre-trained model is more advantageous than leveraging a standard language model. Additionally, we find that math-specific training not only boosts mathematical proficiency but also leads to notable gains in general reasoning benchmarks such as MMLU \citep{mmlu} and BBH \citep{bbh}, highlighting improvements beyond purely mathematical domains.

Following the pre-training phase, we subject DeepSeekMath-Base to mathematical instruction fine-tuning utilizing chain-of-thought \citep{cot}, program-of-thought \citep{pot,pal}, and tool-augmented reasoning datasets \citep{tora}. The refined model, DeepSeekMath-Instruct 7B, surpasses all other models in the 7B parameter class and matches the performance levels of 70B-sized open-source instruction-tuned counterparts.

In addition, we present Group Relative Policy Optimization (GRPO), a reinforcement learning (RL) approach derived from Proximal Policy Optimization (PPO) \citep{schulman2017proximal}. GRPO dispenses with the need for a critic model by inferring baselines from grouped score distributions, greatly reducing computational overhead. When applying GRPO to a subset of English instruction-tuning data, we achieve significant performance enhancements over the robust DeepSeekMath-Instruct, evidenced by both in-domain improvements (GSM8K: 82.9\% $\rightarrow$ 88.2\%, MATH: 46.8\% $\rightarrow$ 51.7\%) and increased out-of-domain results (e.g., CMATH: 84.6\% $\rightarrow$ 88.8\%) during RL training. We further articulate a comprehensive framework that unifies various techniques—including Rejection Sampling Fine-Tuning (RFT) \citep{yuan2023scaling}, Direct Preference Optimization (DPO) \citep{dpo}, PPO, and GRPO—demonstrating that these can all be interpreted as direct or simplified forms of RL. We conduct thorough experimental analyses, comparing aspects such as online versus offline optimization, outcome-oriented versus process-oriented supervision, and single-step versus iterative RL, to gain deeper insights into the factors that underlie this unified framework. Finally, we elucidate the mechanisms by which our RL approach enhances instruction-tuned models and outline promising future directions for further improving RL effectiveness within this unified conceptual landscape.

\subsection{Contributions}
This work introduces advancements in scalable mathematical pre-training, coupled with a systematic investigation into reinforcement learning methodologies.

\textbf{Math Pre-Training at Scale}

\begin{itemize}[topsep=0pt]
    \item Our study demonstrates that Common Crawl, a freely available data source, encompasses substantial mathematical content. Through a rigorous data filtration framework, we curate the DeepSeekMath Corpus—a collection of 120B tokens extracted from web content with an emphasis on mathematics. This dataset is nearly seven times larger than the math data utilized by Minerva \citep{minerva} and approximately nine times greater in scale compared to OpenWebMath \citep{openwebmath}.

    \item We present DeepSeekMath-Base 7B, our foundational pre-trained model, which attains performance levels on par with Minerva 540B \citep{minerva}. These results underscore that parameter count is not the sole determinant of mathematical reasoning proficiency; models with a more modest scale but trained on curated, high-quality data are capable of achieving robust performance.

    \item Insights gained from our math training regimes reveal that introducing code pre-training before mathematical pre-training significantly enhances problem-solving abilities in both tool-based and tool-free settings. This provides empirical evidence relevant to the enduring query: \textit{does exposure to code improve reasoning capabilities?} Our findings affirm this, particularly in the realm of mathematical reasoning.

    \item Contrary to prevailing practices, supplementing training with arXiv papers—a common technique in many mathematics-focused studies—does not yield significant gains on any of the mathematical benchmarks considered in this work.
\end{itemize}

\textbf{Exploration and Analysis of Reinforcement Learning}

\begin{itemize}[topsep=0pt]
    \item We present Group Relative Policy Optimization (GRPO), a resource-efficient and high-performing reinforcement learning algorithm. Distinct from traditional approaches like Proximal Policy Optimization (PPO), GRPO eliminates the use of a critic model and leverages group score baselines, which considerably reduces the computational demands during training.
    \item Our empirical studies reveal that GRPO substantially elevates the capabilities of our instruction-tuned system, DeepSeekMath-Instruct, relying solely on instruction-tuning datasets. Moreover, we document improved out-of-domain generalization throughout the reinforcement learning phase.
    \item We offer an integrated framework for interpreting a range of methods—including RFT, DPO, PPO, and GRPO. In addition, we perform thorough experimental analyses covering various axes such as online versus offline optimization, outcome versus process supervision, and single-turn compared to iterative reinforcement learning, among others, to illuminate the foundational components underpinning this framework.
    \item Building on this unifying framework, we analyze the driving forces behind reinforcement learning’s effectiveness and outline promising research trajectories for enhancing large language models (LLMs) through reinforcement learning.
\end{itemize}

\subsection{Summary of Evaluations and Metrics}

\begin{itemize}[topsep=0pt]
    \item \textbf{English and Chinese Mathematical Reasoning}: Our evaluation encompasses rigorous testing on both English and Chinese datasets spanning mathematical questions from elementary to collegiate levels. The English test suites comprise GSM8K \citep{gsm8k}, MATH \citep{MATH}, SAT \citep{llemma}, OCW Courses \citep{minerva}, and MMLU-STEM \citep{mmlu}, while the Chinese benchmarks feature MGSM-zh \citep{mgsm}, CMATH \citep{wei2023cmath}, Gaokao-MathCloze \citep{agieval}, and Gaokao-MathQA \citep{agieval}. The evaluations scrutinize both the models' capacities for generating independent, step-by-step solutions and their ability to address tasks through Python code.

    On English benchmarks, DeepSeekMath-Base performs on par with the proprietary Minerva 540B \citep{minerva}, outperforming leading open-source foundation models—including Mistral 7B \citep{mistral} and Llemma-34B \citep{llemma}—by a considerable extent, independent of their math pre-training status. Interestingly, DeepSeekMath-Base exhibits even greater prowess on Chinese test suites, which may be attributed to our approach of integrating diverse, high-caliber non-English mathematical data, instead of limiting pre-training to English resources as in previous works \citep{minerva,llemma}. Upon the introduction of instruction-based tuning and reinforcement learning, both DeepSeekMath-Instruct and DeepSeekMath-RL achieve state-of-the-art open-source results, surpassing the 50\% accuracy threshold on the challenging MATH dataset—a first within the open-source landscape.
\end{itemize}

\item \textbf{Formal Mathematics}: DeepSeekMath-Base is assessed on its ability to convert informal statements into formal proofs by leveraging the benchmark task outlined in \citep{dsp_proof}, with evaluations conducted on miniF2F \citep{minif2f} and Isabelle \citep{isabelle} serving as the designated proof assistant. Results show that DeepSeekMath-Base achieves high performance in few-shot autoformalization tasks.

\item \textbf{Natural Language Understanding, Reasoning, and Code}: To thoroughly characterize the general capabilities of the models in understanding, reasoning, and code generation, DeepSeekMath-Base is evaluated using several benchmarks. These include the Massive Multitask Language Understanding (MMLU) suite \citep{mmlu} spanning 57 multiple-choice topics, BIG-Bench Hard (BBH) \citep{bbh} which features 23 tasks emphasizing multi-step reasoning, as well as HumanEval \citep{codex} and MBPP \citep{mbpp}, both of which are standard datasets for assessing coding proficiency. Pre-training on mathematical data leads to improvements in both language comprehension and reasoning skills.

\section{Math Pre-Training}

\subsection{Data Collection and Decontamination} This section details the procedure used to assemble the DeepSeekMath Corpus by extracting data from Common Crawl. As shown in Figure \ref{fig:data_collect}, our method adopts an iterative framework designed to efficiently build a large-scale math corpus from Common Crawl, initiating with a high-quality seed collection of mathematical texts. This methodology can similarly be extended to other areas such as computer code.

Initially, OpenWebMath \citep{openwebmath}—a curated set of high-caliber mathematical texts from the web—is selected as the foundational seed corpus. With this dataset, a fastText model \citep{joulin2016fasttext} is trained to retrieve additional web documents resembling OpenWebMath. The process involves randomly drawing 500,000 samples from the seed set as positive training instances and pairing them with 500,000 negative samples from Common Crawl. Training leverages an open-source toolkit\footnote{\url{https://fasttext.cc}}, employing a vector size of 256, a learning rate of 0.1, a maximum word n-gram length of 3, a minimum word frequency of 3, and 3 training epochs. To manage the large scale of Common Crawl, URL-based deduplication and near-deduplication are applied, reducing the dataset to 40 billion HTML pages. Using the trained fastText classifier, mathematical web pages are extracted from this deduplicated set. To ensure high-quality mathematical content, the harvested pages are sorted according to their fastText model scores, retaining only those with the highest rankings. Several pre-training experiments are conducted using the top 40B, 80B, 120B, and 160B tokens to determine optimal data volume. For the initial iteration, the decision is made to keep the leading 40B tokens.

Following the initial round of data collection, a significant number of mathematical web pages remain missing. This shortfall arises largely because the set of positive examples used to train the fastText model lacks adequate diversity. To address this, we seek out additional math-focused web domains to supplement the seed corpus and thereby refine the fastText classifier. Our procedure involves segmenting the entire Common Crawl dataset into non-overlapping domains, where a domain encompasses all pages sharing a common base URL. For each domain, we compute the proportion of pages captured during the first pass. Domains with more than 10\% of their pages collected are categorized as mathematics-oriented (such as \url{mathoverflow.net}). 

Next, we conduct a manual review to annotate URLs within these domains that are linked to mathematical content (for example, \url{mathoverflow.net/questions}). Any web pages under these annotated URLs that have not yet been collected are incorporated into our growing seed set. This strategy allows us to expand our pool of positive training samples and, in turn, enhances the capability of the fastText model to identify and retrieve additional mathematical web pages in subsequent rounds. After four cycles of collection, we assemble a dataset comprising 35.5M mathematical web pages and amounting to 120B tokens. In the fourth round, we observe that nearly 98\% of the material had already been gathered by the third, prompting us to discontinue further collection.

To ensure the math corpus does not contain evaluation set leakage, we adopt the methodology in \cite{deepseek-coder} to exclude web pages that contain content from English mathematical benchmarks, including GSM8K \citep{gsm8k} and MATH \citep{MATH}, as well as Chinese benchmarks like CMATH \citep{wei2023cmath} and AGIEval \citep{agieval}. Our exclusion rules are as follows: if any text snippet features a 10-gram sequence that is an exact match for a substring in any benchmark, that snippet is removed from the training set. For benchmark strings shorter than 10 grams but at least 3 grams in length, we apply exact matching to identify and remove any contaminated web content.

\subsection{Validating the Quality of the DeepSeekMath~Corpus}

We conduct pre-training trials to compare the DeepSeekMath~Corpus against contemporary, large-scale math training corpora:

\begin{itemize}[topsep=0pt]
    \item \textbf{MathPile} \citep{mathpile}: An 8.9B-token composite dataset derived from sources including textbooks, Wikipedia, ProofWiki, CommonCrawl, StackExchange, and arXiv—over 85\% of which originates from arXiv;
    \item \textbf{OpenWebMath} \citep{openwebmath}: A 13.6B-token corpus filtered from CommonCrawl specifically for mathematical material;
    \item \textbf{Proof-Pile-2} \citep{llemma}: A mathematics-focused collection that comprises OpenWebMath, AlgebraicStack (providing 10.3B tokens of mathematical code), and arXiv manuscripts (totaling 28.0B tokens). In experiments involving Proof-Pile-2, we align with \cite{llemma} and maintain an arXiv:Web:Code distribution ratio of 2:4:1.
\end{itemize}

\subsubsection{Training Setting}
\label{sec:quality-policy}
We fine-tune a general-purpose language model consisting of 1.3B parameters on mathematical datasets, using the architecture underlying DeepSeek LLMs \citep{deepseek-llm}, hereafter referred to as DeepSeek-LLM 1.3B. Each mathematical dataset serves as the basis for an independent training run, each spanning 150B tokens. The entirety of our training workflow utilizes the resource-efficient HAI-LLM framework \citep{haillm}. Adhering to the established protocols for DeepSeek LLMs, we leverage the AdamW optimizer \citep{adamW} configured with $\beta_1=0.9$, $\beta_2=0.95$, and a $\mathrm{weight\_decay}$ of 0.1. The learning rate follows a staged schedule: it linearly increases to its maximum value over 2,000 warmup steps, drops to 31.6\% of the maximum by the 80\% mark of total training, and then further decays to 10.0\% of its peak after 90\% of the total steps. The peak learning rate is set at 5.3e-4. We adopt a batch size comprising 4M tokens and a maximum sequence length of 4K.

\subsubsection{Evaluation Results}

\textbf{The DeepSeekMath~Corpus demonstrates superior quality, extensive multilingual coverage, and unprecedented scale compared to prior work.}

\begin{itemize}[topsep=0pt]
    \item \textbf{High-quality}:
    Model performance is assessed using eight mathematical benchmarks, employing few-shot chain-of-thought prompting \cite{cot}. Table \ref{tab:corpora_comparison} highlights that the model trained on the DeepSeekMath~Corpus substantially outperforms those trained on other datasets. As further illustrated in Figure \ref{fig:corpora_comparisons}, the DeepSeekMath-trained model surpasses the model trained on Proof-Pile-2 after processing 50B tokens (equivalent to one complete epoch on Proof-Pile-2), underscoring the superior average quality of the DeepSeekMath~Corpus.
    \item \textbf{Multilingual}:
    The DeepSeekMath~Corpus comprises material in various languages, with English and Chinese constituting the dominant language groups. According to Table \ref{tab:corpora_comparison}, models trained on the DeepSeekMath~Corpus exhibit enhanced mathematical reasoning in both English and Chinese. By contrast, existing mathematical datasets are mainly English-focused, providing limited improvements in Chinese reasoning or, in some cases, even impairing performance.
    \item \textbf{Large-scale}:
    The DeepSeekMath~Corpus vastly exceeds current mathematical corpora in magnitude. As shown in Figure \ref{fig:corpora_comparisons}, DeepSeek-LLM 1.3B trained on this corpus achieves a much sharper learning trajectory and sustained performance gains. Conversely, baseline datasets are significantly smaller and have already undergone multiple passes within a training cycle, causing the corresponding model improvements to quickly level off.
\end{itemize}

\subsection{Training and Evaluation of DeepSeekMath-Base 7B}

This section details the development of DeepSeekMath-Base 7B, a foundational model characterized by robust reasoning skills, with particular emphasis on mathematical competence. DeepSeekMath-Base 7B is initialized using DeepSeek-Coder-Base-v1.5 7B \citep{deepseek-coder} and undergoes training on 500B tokens. The dataset composition is as follows: 56\% from the DeepSeekMath Corpus, 4\% sourced from AlgebraicStack, 10\% drawn from arXiv, 20\% consisting of Github code, and the remaining 10\% comprising natural language material from Common Crawl in both English and Chinese. The training procedure largely mirrors that described in Section \ref{sec:quality-policy}, with the exception that the maximum learning rate is adjusted to 4.2e-4 and the batch size is set at 10 million tokens.

A thorough evaluation of DeepSeekMath-Base 7B is performed to measure its mathematical proficiency. The evaluation considers its aptitude for generating comprehensive, self-contained mathematical solutions independent of external tools, its effectiveness at solving problems utilizing appropriate tools, and its capability in formal theorem proving. In addition to its mathematical aptitude, we provide an overview of the base model’s performance in broader domains, including natural language understanding, logical reasoning, and programming tasks.

\paragraph{Mathematical Problem Solving via Stepwise Reasoning}

To gauge DeepSeekMath-Base's aptitude for mathematical problem solving, we employ few-shot chain-of-thought prompting \citep{cot} across eight diverse benchmarks in both English and Chinese. The benchmarks include quantitative reasoning datasets (such as GSM8K \citep{gsm8k}, MATH \citep{MATH}, and CMATH \citep{wei2023cmath}) and multiple-choice formats (such as MMLU-STEM \citep{mmlu} and Gaokao-MathQA \citep{agieval}), spanning a range of mathematical topics from basic to collegiate levels.

As detailed in Table \ref{tab:base_math_cot}, DeepSeekMath-Base 7B achieves top results across all eight evaluated benchmarks compared to other open-source base models, including both the extensively utilized Mistral 7B \citep{mistral} and the recently introduced Llemma 34B \citep{llemma}, which received mathematical training on Proof-Pile-2 \citep{llemma}. Particularly remarkable is DeepSeekMath-Base’s performance on the advanced MATH dataset, where it surpasses prior open-source base models by more than 10 percentage points, and even exceeds the capabilities of Minerva 540B \citep{minerva}—a proprietary model with 77 times the parameter count, derived from PaLM \citep{palm} and further refined on mathematical corpora.

\paragraph{Mathematical Problem Solving with Tool Use} We assess mathematical reasoning that incorporates tool assistance on GSM8K and MATH by employing few-shot program-of-thought prompting strategies \citep{pot,pal}. In this setup, models address each task by constructing Python scripts, optionally making use of packages such as \textit{math} and \textit{sympy} to handle complex calculations. The answer is determined by executing the generated code. According to Table \ref{tab:base_math_pot_proof}, DeepSeekMath-Base 7B sets a new performance benchmark, exceeding the results achieved by the former leader, Llemma 34B.

\paragraph{Formal Mathematics}

Automated formalization of proofs contributes to the rigor and dependability of mathematical argumentation while streamlining the proof process—a topic attracting growing interest. We benchmark DeepSeekMath-Base 7B on the informal-to-formal proof generation task presented in \citep{dsp_proof}, which involves producing a formal proof given an informal proposition, its formal expression, and an informal argument. Evaluations are conducted on the miniF2F dataset \citep{minif2f}, which tests formal problem-solving abilities at the Olympiad level; the model generates formal proofs in Isabelle for each prompt using a few-shot prompting approach. As recommended by \cite{dsp_proof}, we let models create proof outlines and then apply the automated proof assistant Sledgehammer \citep{sledgehammer} to complete the detailed steps. Table \ref{tab:base_math_pot_proof} indicates that DeepSeekMath-Base 7B delivers notable outcomes in the autoformalization of proofs.

\paragraph{Natural Language Understanding, Reasoning, and Code}

For natural language understanding, we use MMLU \citep{mmlu}; for evaluating reasoning, we consider BBH \citep{bbh}; and for coding performance, we employ HumanEval \citep{codex} and MBPP \citep{mbpp}. Table \ref{tab:understanding_reasoning_code} reveals that DeepSeekMath-Base 7B demonstrates marked improvements on MMLU and BBH compared to its predecessor, DeepSeek-Coder-Base-v1.5 \citep{deepseek-coder}, highlighting the benefits of targeted math training on both comprehension and reasoning abilities. Moreover, the addition of code-related tokens in further training enables DeepSeekMath-Base 7B to sustain the coding proficiency of DeepSeek-Coder-Base-v1.5 across both programming benchmarks. In summary, DeepSeekMath-Base 7B achieves significant gains relative to the general-purpose Mistral 7B \citep{mistral} model on all three evaluated reasoning and coding tasks.

\section{Supervised Fine-Tuning}

\subsection{SFT Data Curation}

We assemble an instruction-tuning dataset for mathematics that encompasses both English and Chinese problems, sourced from a wide range of mathematical disciplines and complexity levels. Each problem is paired with corresponding solutions that adopt chain-of-thought (CoT) \citep{cot}, program-of-thought (PoT) \citep{pot,pal}, and tool-integrated reasoning styles \citep{tora}. Our curated dataset consists of a total of 776K training instances.

\begin{itemize}[topsep=0pt]
    \item \textbf{English mathematical datasets}: Our English collection features annotated versions of GSM8K and MATH, where solutions are constructed to incorporate tool-based reasoning. In addition, we include selected entries from MathInstruct \citep{MathInstruct} and the training set from Lila-OOD \citep{lila}, both presenting solutions via CoT or PoT. The resulting corpus reflects a broad array of mathematical subjects, such as algebra, probability, number theory, calculus, and geometry.
    \item \textbf{Chinese mathematical datasets}: We aggregate Chinese K-12 mathematics questions drawn from 76 different sub-domains, for instance, linear equations, providing both chain-of-thought and tool-integrated solution annotations.
\end{itemize}

\subsection{Training and Evaluating DeepSeekMath-Instruct 7B}

This section details the development of DeepSeekMath-Instruct 7B, a model subjected to instruction-tuning in mathematics, beginning with DeepSeekMath-Base as its foundation. Training samples are concatenated in random order until the cumulative context size reaches 4K tokens. The training process spans 500 iterations, utilizing a batch size of 256 and maintaining a fixed learning rate of 5e-5.

For evaluation, we assess the mathematical reasoning abilities of our models—both in standard settings and in conjunction with external tools—on four quantitative reasoning benchmarks in English and Chinese. To contextualize performance, we compare our model against the current leading alternatives:

\begin{itemize}[topsep=0pt]
    \item \textbf{Closed-source models} under comparison encompass: (1) the GPT series, with GPT-4 \citep{gpt4} and GPT-4 Code Interpreter~\footnote{\url{https://openai.com/blog/chatgpt-plugins\#\#code-interpreter}} being among the most advanced; (2) Gemini Ultra and Pro \citep{gemini}; (3) Inflection-2 \citep{inflection-2}; (4) Grok-1~\footnote{\url{https://x.ai/model-card}}; as well as several recently introduced models from Chinese enterprises, such as (5) Baichuan-3~\footnote{\url{https://www.baichuan-ai.com}}, and (6) the latest member of the GLM family, GLM-4~\footnote{\url{https://open.bigmodel.cn/dev/api\#glm-4}} \citep{glm}.
\end{itemize}

These models are intended for broad applications and most have been subjected to various alignment techniques.

\item \textbf{Open-source models} comprise: general-purpose systems such as (1) DeepSeek-LLM-Chat 67B \citep{deepseek-llm}, (2) Qwen 72B \citep{qwen}, (3) SeaLLM-v2 7B \citep{seallm}, and (4) ChatGLM3 6B \citep{chatglm3}; as well as models optimized for mathematical tasks, including (5) InternLM2-Math 20B\footnote{\url{https://github.com/InternLM/InternLM-Math}}—an extension of InternLM2 with additional mathematics training and subsequent instruction fine-tuning, (6) Math-Shepherd-Mistral 7B, which incorporates PPO training \citep{schulman2017proximal} using a process-supervised reward model on top of Mistral 7B \citep{mistral}, (7) the WizardMath series \citep{wizardmath}, which leverages evolve-instruct (a variant of instruction tuning with AI-generated instructions) and PPO training to enhance Mistral 7B and Llama-2 70B \citep{llama2} for mathematical reasoning, mainly using GSM8K and MATH datasets for training, (8) MetaMath 70B \citep{metamath}, consisting of Llama-2 70B fine-tuned with expanded GSM8K and MATH data, (9) ToRA 34B \cite{tora}, representing CodeLlama 34B adapted for tool-assisted mathematical reasoning, and (10) MAmmoTH 70B \citep{MathInstruct}, where Llama-2 70B is refined via instruction tuning using MathInstruct.

As displayed in Table \ref{tab:sft_rl_math}, when tool integration is not permitted during evaluation, DeepSeekMath-Instruct 7B delivers robust stepwise reasoning. Of particular note, on the challenging MATH benchmark, our model achieves an absolute performance gain of at least 9\% over all other open-source systems and most commercial offerings (such as Inflection-2 and Gemini Pro), even outperforming models with significantly more parameters (e.g., Qwen 72B) or those extensively trained with math-centric reinforcement learning strategies (e.g., WizardMath-v1.1 7B). While DeepSeekMath-Instruct’s results are on par with proprietary Chinese models like GLM-4 and Baichuan-3 for MATH, it continues to trail behind leading solutions like GPT-4 and Gemini Ultra.

In settings that allow a hybrid approach—combining natural language and code-based tool use for solving problems—DeepSeekMath-Instruct 7B achieves nearly 60\% accuracy on MATH, outstripping all available open-source models. On the other test sets, our model remains competitive with DeepSeek-LLM-Chat 67B, the previous leading system, despite being only a tenth its size.

\section{Reinforcement Learning}

\subsection{Group Relative Policy Optimization} Reinforcement learning (RL) has established itself as a successful strategy for further boosting the mathematical reasoning performance of LLMs after Supervised Fine-Tuning (SFT) \citep{wang2023math, wizardmath}. We now present our novel and efficient RL approach, Group Relative Policy Optimization (GRPO).

\subsubsection{From PPO to GRPO} Proximal Policy Optimization (PPO) \citep{schulman2017proximal} stands out as a widely adopted actor-critic RL method for fine-tuning LLMs during the RL phase \citep{ouyang2022training}. It functions by maximizing the following surrogate objective function:

\begin{equation}
\footnotesize
    \mathcal{J}_{PPO}(\theta) = \mathbb{E}{[q \sim P(Q), o \sim \pi_{\theta_{old}}(O|q)]} \frac{1}{|o|} \sum_{t=1}^{|o|} \min \left[ \frac{\pi_\theta(o_{t} | q, o_{<t})}{\pi_{\theta_{old}}(o_{t} | q, o_{<t})} A_{t}, \text{clip} \left( \frac{\pi_\theta(o_{t} | q, o_{<t})}{\pi_{\theta_{old}}(o_{t} | q, o_{<t})}, 1 - \varepsilon, 1 + \varepsilon \right)  A_{t} \right

Here, $\pi_{\theta}$ denotes the current policy model, while $\pi_{\theta_{old}}$ stands for the previous version of the policy. The variables $q$ and $o$ represent questions and outputs, respectively, both sampled either from the dataset of questions or from the output distribution of the previous policy $\pi_{\theta_{old}}$. The hyper-parameter $\varepsilon$ is employed for clipping in PPO, serving to enhance the stability of the training procedure. The term $A_t$ indicates the advantage, calculated using Generalized Advantage Estimation (GAE) \citep{gae}, based on rewards $\{r_{\ge t}\}$ and a parameterized value function $V_{\psi}$. Consequently, PPO necessitates the concurrent training of a value function together with the policy. To prevent excessive optimization towards the reward model, a conventional technique involves including a per-token Kullback-Leibler (KL) divergence penalty—relative to a reference model—within the reward for each token, as introduced by \citep{ouyang2022training}:

\begin{equation}
    r_{t} = r_\varphi(q, o_{\le t}) - \beta \log\frac{\pi_{\theta}(o_{t}|q, o_{<t})}{\pi_{ref}(o_{t}|q, o_{<t})},
\label{eq:PPO-reward}
\end{equation}

Here, $r_\varphi$ signifies the reward model, and $\pi_{ref}$ designates the reference policy (typically the initial SFT model). The hyper-parameter $\beta$ modulates the KL penalty's strength.

Because the value function in PPO is often implemented as a neural network with complexity similar to the policy model itself, this imposes significant demands on both memory and computational resources. Moreover, throughout RL optimization, the value function provides a baseline in the computation of the advantage to control variance. Within the large language model (LLM) framework, reward models usually assign nonzero scores only to the final token, which can make accurate, token-level value function learning difficult. To overcome these issues, as illustrated in Figure \ref{fig:grpo}, we introduce Group Relative Policy Optimization (GRPO). Unlike PPO, GRPO dispenses with the value function approximation; instead, it utilizes the mean reward from multiple sampled outputs—each generated for the same prompt—as a baseline. Specifically, given any question $q$, GRPO samples a collection of outputs $\{o_1, o_2, \cdots, o_G\}$ using the previous policy $\pi_{\theta_{old}}$, and optimizes the new policy by maximizing:

\begin{equation}
\footnotesize
\begin{split}
    \mathcal{J}_{GRPO}(\theta) &= \mathbb{E}{[q \sim P(Q), \{o_i\}_{i=1}^G \sim \pi_{\theta_{old}}(O|q)]}  \\
    & \frac{1}{G}\sum_{i=1}^G\frac{1}{|o_i|} \sum_{t=1}^{|o_i|} \left\{ \min \left[ \frac{\pi_\theta(o_{i,t} | q, o_{i,<t})}{\pi_{\theta_{old}}(o_{i,t} | q, o_{i,<t})} \hat{A}_{i,t}, \text{clip} \left( \frac{\pi_\theta(o_{i,t} | q, o_{i,<t})}{\pi_{\theta_{old}}(o_{i,t} | q, o_{i,<t})}, 1 - \varepsilon, 1 + \varepsilon \right)  \hat{A}_{i,t} \right] - \beta \mathbb{D}_{KL}\left[\pi_{\theta} || \pi_{ref}\right]\right\} ,
\end{split}
\label{eq:GRPO-obj}
\end{equation}

where the hyper-parameters $\varepsilon$ and $\beta$ remain, and $\hat{A}_{i,t}$ denotes the group-relative advantage, computed from rewards among outputs within the same group, as will be explained further in subsequent subsections. This strategy for advantage estimation directly mirrors the comparative structure of reward model datasets, where rewards are based on comparisons among answers to identical prompts. It should also be emphasized that, rather than integrating the KL penalty into the reward, GRPO directly incorporates the KL divergence between the new and reference policies into the loss. This design choice simplifies the computation of $\hat{A}_{i,t}$. Unlike the approach in (\ref{eq:PPO-reward}), GRPO uses

\begin{algorithm}[t]
  \small
  \caption{Iterative Group Relative Policy Optimization}
  \textbf{Input} initial policy model $\pi_{\theta_{\text{init}}}$; reward models $r_\varphi$; task prompts $\mathcal{D}$; 
  hyperparameters $\varepsilon$, $\beta$, $\mu$
  \begin{algorithmic}[1]
    \State Initialize policy model $\pi_\theta \leftarrow \pi_{\theta_{\text{init}}}$
    \For{iteration = 1, \dots, I}
       \State Assign current policy as reference model: $\pi_{ref} \leftarrow \pi_{\theta}$
      \For{step = 1, \dots, M}
      \State Select a batch $\mathcal{D}_b$ sampled from $\mathcal{D}$
      \State Store previous policy: $\pi_{\theta_{old}} \leftarrow \pi_{\theta}$ 
      \State For every question $q \in \mathcal{D}_b$, sample $G$ responses $\{o_i\}_{i=1}^G$ using $\pi_{\theta_{old}} (\cdot \mid q)$
      \State Calculate rewards $\{r_i\}_{i=1}^{G}$ for each output $o_i$ by evaluating $r_{\varphi}$
      \State For each token $t$ in $o_i$, compute $\hat{A}_{i,t}$ using group relative advantage estimation
      \For{GRPO iteration = 1, \dots, $\mu$}
        \State Update $\pi_{\theta}$ to increase the GRPO objective as specified in Equation \ref{eq:GC-GRPO}
      \EndFor
    \EndFor 
    \State Refine $r_\varphi$ via ongoing training with a replay buffer.
    \EndFor 
  \end{algorithmic}
  \textbf{Output} $\pi_\theta$
  \label{alg:iter-grpo}
\end{algorithm}

\subsubsection{Outcome Supervision RL with GRPO} More formally, for each prompt $q$, a set of $G$ candidate outputs $\{o_1, o_2, \cdots, o_G\}$ are generated using the previous policy $\pi_{\theta_{old}}$. Each output is then evaluated by a reward model, producing a reward vector $\mathbf{r} = \{r_1, r_2, \cdots, r_G\}$. The set of rewards is then standardized by subtracting the mean and dividing by the standard deviation within the group. In the outcome supervision setting, the resulting normalized reward for output $o_i$ is assigned to every token position as the advantage: that is, for every $t$ in $o_i$, $\hat{A}_{i, t} = \widetilde{r}_i = \frac{r_i- {\rm mean}(\mathbf{r})}{{\rm std}(\mathbf{r})}$. The policy is then updated to maximize the objective indicated in equation (\ref{eq:GRPO-obj}).

\subsubsection{Process Supervision RL with GRPO} While outcome supervision only supplies a final reward at the end of each output, it may not provide sufficiently granular feedback for the policy, especially on tasks requiring stepwise mathematical reasoning. In line with \cite{wang2023math}, we incorporate process supervision, where feedback is issued after every reasoning stage. Formally, after generating $G$ sample outputs $\{o_1, o_2, \cdots, o_G\}$ to a given question $q$, a process-level reward model evaluates each intermediate step, producing: $\mathbf{R} = \{ \{r_1^{index(1)},\cdots,r_1^{index(K_1)}\}, \cdots,  \{r_G^{index(1)},\cdots,r_G^{index(K_G)}\} \}$, where $index(j)$ marks the terminal token of the $j$-th reasoning step, and $K_i$ is the number of steps for output $o_i$. Rewards for all steps are normalized via group mean and standard deviation, i.e., $\widetilde{r}_i^{index(j)} = \frac{r_i^{index(j)} - {\rm mean(\mathbf{R})}}{{\rm std(\mathbf{R})}}$. The advantage $\hat{A}_{i, t}$ for token $t$ is then the cumulative sum of subsequent normalized rewards for that output, specifically: $\hat{A}_{i, t} = \sum_{index(j) \ge t} \widetilde{r}_i

\subsubsection{Iterative RL with GRPO} During reinforcement learning, reliance on an outdated reward model may fail to provide adequate supervision as the policy model evolves. To address this, we investigate an iterative RL strategy employing GRPO. As depicted in Algorithm \ref{alg:iter-grpo}, iterative GRPO constructs updated training datasets for the reward model by sampling new outputs from the current policy model, then continually updates the reward model via a replay method that maintains 10\% of data from previous iterations. Subsequently, the reference model is updated to mirror the latest policy model, and training resumes with the policy model under the supervision of the newly refined reward model.

\subsection{Training and Evaluating DeepSeekMath-RL}

Our reinforcement learning experiments are based on DeepSeekMath-Instruct 7B. The RL dataset comprises chain-of-thought questions sampled from GSM8K and MATH within the SFT data, totaling roughly 144K instances. To analyze how RL influences performance on benchmarks not present in the RL dataset, we exclude all other SFT questions. Following the procedure in \citep{wang2023math}, the reward model's training set is constructed accordingly. The initial reward model is trained with DeepSeekMath-Base 7B and a learning rate of 2e-5. For policy model optimization within the GRPO framework, we adopt a learning rate of 1e-6 and set the KL divergence coefficient to 0.04. Each question prompts the generation of $64$ distinct outputs. The maximum output length is limited to 1024 tokens, and training is performed with a batch size of 1024. Policy model updates occur only once per exploration phase. For evaluation, DeepSeekMath-RL 7B is assessed on the same suite of benchmarks as DeepSeekMath-Instruct 7B. In this context, tasks drawn from GSM8K and MATH with chain-of-thought solutions are categorized as in-domain, while all remaining benchmarks are considered out-of-domain.

Table \ref{tab:sft_rl_math} presents the comparative results for both open- and closed-source models, considering their performance on chain-of-thought and tool-integrated reasoning tasks in English and Chinese evaluations. Our observations reveal the following: (1) DeepSeekMath-RL 7B achieves accuracies of 88.2\% on GSM8K and 51.7\% on MATH when employing chain-of-thought strategies. These scores not only set a new benchmark among all open-source models within the 7B to 70B scale but also outperform most leading closed-source models. (2) Notably, DeepSeekMath-RL 7B's training is limited solely to chain-of-thought-based instruction tuning data from GSM8K and MATH, beginning from DeepSeekMath-Instruct 7B. Despite the narrowly focused training data, DeepSeekMath-RL 7B consistently exceeds DeepSeekMath-Instruct 7B across all measured benchmarks, underscoring the value of reinforcement learning.

\section{Discussion}
This section outlines our key insights gained from both pre-training and reinforcement learning studies.

\subsection{Lessons Learnt in Pre-Training}

We begin by discussing pre-training experiences. Except where explicitly indicated, the experimental procedures correspond to those in Section \ref{sec:quality-policy}. It should also be highlighted that references to the DeepSeekMath Corpus in this context pertain to an 89B-token dataset collected during the project’s second data gathering iteration.

\subsubsection{Code Training Benefits Mathematical Reasoning}

The idea that code training can enhance reasoning abilities is frequently posited, but concrete evidence remains limited. Our investigation sheds light on this claim within the mathematical reasoning context: code-oriented training enhances a model’s mathematical reasoning skills, irrespective of tool assistance.

To explore the influence of code training on mathematical reasoning, we analyzed performance using two-stage and one-stage training paradigms:

\textbf{Two-Stage Training}

\begin{itemize}[topsep=0pt]
    \item \textbf{Code Training for 400B Tokens $\rightarrow$ Math Training for 150B Tokens}: We first train DeepSeek-LLM 1.3B on a corpus of 400B code tokens, then subsequently expose it to 150B tokens focused on mathematics;
    \item \textbf{General Training for 400B Tokens $\rightarrow$ Math Training for 150B Tokens}: For comparison, we substitute code tokens with general tokens (sampled from DeepSeek-AI's comprehensive general corpus) during the initial 400B token phase, aiming to evaluate whether code-centric pre-training offers distinct advantages for developing mathematical reasoning over a more generalized approach.
\end{itemize}

\textbf{One-Stage Training}

\begin{itemize}[topsep=0pt]
    \item \textbf{Math Training for 150B Tokens}: DeepSeek-LLM 1.3B is exclusively trained on 150B tokens dedicated to mathematics;
    \item \textbf{Training on a mixture of 400B Code Tokens and 150B Math Tokens}: Given that a sequential math training phase after code training reduces performance on coding tasks, we also test simultaneous exposure by combining 400B code tokens with 150B math tokens in a single training stage, with the aim of determining if this joint regimen fosters mathematical reasoning gains while avoiding catastrophic forgetting.
\end{itemize}

\paragraph{Results}
Performance metrics under these training strategies are provided in Table \ref{tab:code-math} and Table \ref{tab:code-math-general-eval}.

Code-focused pre-training enhances the model’s program-assisted mathematical reasoning abilities in both two-stage and one-stage regimens. Table \ref{tab:code-math} illustrates that two-stage training, starting with code and followed by mathematics, markedly increases success rates on tasks like GSM8K and MATH (evaluated in a Python-execution setting). An additional round of math-specific training in the second phase further amplifies performance. Notably, in the one-stage mixed regime, the integration of code and math tokens diminishes the negative impacts of catastrophic forgetting encountered in sequential training, and leads to improvements in both code (Table \ref{tab:code-math-general-eval}) and program-augmented mathematical problem solving (Table \ref{tab:code-math}).

Moreover, code pre-training confers improvements on pure mathematical reasoning, even without recourse to tool-use. The first phase of code training in the two-stage pipeline yields a notable, though moderate, boost on direct reasoning tasks, and paves the way for more effective subsequent math-centric learning—eventually achieving optimal results. In contrast, in the one-stage blended-token setup, mathematical reasoning performance without tool support is negatively affected. A plausible explanation is that DeepSeek-LLM 1.3B’s limited model capacity constrains its ability to jointly assimilate and benefit from both code and math modalities when exposed to both simultaneously.

\subsubsection{ArXiv Papers Exhibit Limited Effectiveness in Enhancing Mathematical Reasoning}

Although arXiv papers are frequently incorporated into mathematical pre-training datasets \citep{minerva,gpt-f,llemma,mathpile}, there is a scarcity of comprehensive investigations examining their contribution to mathematical reasoning abilities. Contrary to prevailing assumptions, our experimental findings suggest that arXiv papers do not significantly bolster mathematical reasoning. We assess this across models of varying scales, such as DeepSeek-LLM 1.3B and DeepSeek-Coder-Base-v1.5 7B \citep{deepseek-coder}, using distinct arXiv datasets processed via different methodologies:

\begin{itemize}[topsep=0pt]
    \item \textbf{MathPile} \citep{mathpile}: a curated corpus totaling 8.9B tokens, compiled using systematic cleaning and filtering protocols, of which over 85\% consists of scientific arXiv papers;
    \item \textbf{ArXiv-RedPajama} \citep{redpajama}: a dataset comprising all arXiv LaTeX source files, with non-essential content such as preambles, comments, macros, and bibliographies removed, amounting to 28.0B tokens.
\end{itemize}

In these studies, we independently train DeepSeek-LLM 1.3B with 150B tokens and DeepSeek-Coder-Base-v1.5 7B with 40B tokens for each arXiv-derived dataset. Results consistently reveal that relying solely on arXiv papers yields little to no enhancement—sometimes even a decline—in mathematical reasoning performance, as assessed by a spectrum of benchmarks spanning different difficulty levels. The evaluation includes quantitative reasoning sets such as GSM8K and MATH (see Table \ref{tab:arxiv-cot}), STEM-oriented multiple-choice tasks like MMLU-STEM (Table \ref{tab:arxiv-cot}), and formal mathematics evaluations including miniF2F (Table \ref{tab:arxiv_atp}).

Nevertheless, this observation is subject to certain caveats. Several aspects remain unexplored, including:

\begin{itemize}[topsep=0pt]
    \item The potential effects of arXiv-derived tokens on specialized mathematical tasks not addressed in this work, such as the informalization of theorems—that is, transforming formal proofs or statements into their informal counterparts;
    \item The interactions between arXiv content and other categories of training data;
    \item Possible advantages conferred by arXiv pre-training at much larger model scales.
\end{itemize}

Consequently, additional in-depth analyses are warranted, and we reserve this for future investigation.

\subsection{Insights on Reinforcement Learning}
\subsubsection{Toward a Unified Framework} This section introduces a unifying framework for conceptualizing a variety of training approaches, including SFT, RFT, DPO, PPO, and GRPO, and presents empirical studies to examine influential elements within this paradigm. Broadly, the parameter gradient for a given training method is formulated as:

\begin{equation}
    \nabla_{\theta}\mathcal{J}_{\textcolor{red}{\mathcal{A}}}(\theta) = \mathbb{E}[\underbrace{(q,o) \sim \textcolor{red}{\mathcal{D}}}_{Data \ Source}]\left( \frac{1}{|o|} \sum_{t=1}^{|o|}  \underbrace{GC_{{\mathcal{A}}}(q, o, t, \textcolor{red}{\pi_{{rf}}})}_{Gradient \ Coefficient}  \nabla_{\theta}\log \pi_{\theta}(o_t | q, o_{<t})\right).
\label{eq:objective}
\end{equation}

This framework is characterized by three principal elements: 1) the \textit{Data Source} $\mathcal{D}$, which supplies the dataset used for model training; 2) the \textit{Reward Function} $\pi_{{rf}}$, responsible for generating the reward signals guiding learning; and 3) the \textit{Algorithm} $\mathcal{A}$, which integrates the training examples and reward signals to yield the gradient coefficient $GC$, thus modulating the strength and direction of updates (either penalizing or rewarding certain behaviors). We review a selection of prominent methods under this common structure:

\begin{itemize}
    \item \textbf{Supervised Fine-tuning (SFT)}: In SFT, a pretrained model is further optimized using human-curated data tailored for fine-tuning.
    \item \textbf{Rejection Sampling Fine-tuning (RFT)}: RFT involves an additional fine-tuning step for the SFT model by utilizing outputs generated in response to SFT prompts, retaining only those responses deemed correct for subsequent training.
    \item \textbf{Direct Preference Optimization (DPO)}: DPO extends SFT by performing additional fine-tuning, leveraging both original and newly generated outputs together with a pairwise preference loss.
    \item \textbf{Online Rejection Sampling Fine-tuning (Online RFT)}: Diverging from standard RFT, Online RFT uses the SFT-initialized model as the starting policy and dynamically updates the policy through fine-tuning on samples produced by the current version of the policy itself.
    \item \textbf{PPO/GRPO}: In these approaches, the model is initialized using SFT parameters and then improved further using real-time sampling and reinforcement signals from the evolving policy model.
\end{itemize}

The main elements of these approaches are outlined in Table \ref{tab:data_policy}. For an in-depth explanation of the derivation, please consult Appendix \ref{app:analysis-rl}.

\paragraph{Analysis of Data Source} We categorize data sources into two groups: online and offline sampling. Online sampling implies that training data is generated from real-time outputs of the policy model during exploration, whereas offline sampling utilizes data originating from the initial SFT model's sampling. RFT and DPO operate under the offline paradigm, while Online RFT and GRPO employ an online data acquisition strategy.

Referencing Figure \ref{fig:rl-analysis}, we observe that Online RFT achieves notably superior performance over RFT on two evaluation benchmarks. Early in the training process, Online RFT and RFT show similar results; however, as training advances, Online RFT establishes a clear advantage, underscoring the benefits of online learning. This outcome aligns with expectations: at the start, the behavior of the actor and SFT model is largely indistinguishable, so sampled data yields minimal differences. As training proceeds, samples from the evolving actor increasingly diverge, and dynamic data collection becomes significantly more advantageous.

\paragraph{Analysis of Gradient Coefficient} The transformation of the reward signal into a gradient coefficient is essential for updating model parameters. In our setup, the reward function is either based on explicit `Rule’ criteria or on an auxiliary `Model’ trained to evaluate responses. Rule-based rewards assess correctness directly, while model-based rewards use a learned reward model, itself trained on rule-based evaluations. Equations \ref{eq:GC-RFT} and \ref{eq:GC-GRPO} illustrate a fundamental distinction: GRPO modulates its gradient coefficient depending on reward values output by the reward model, thereby selectively enhancing or suppressing responses according to their reward magnitude. In contrast, Online RFT does not differentiate in this way—it fails to penalize incorrect responses, instead reinforcing all correct answers uniformly without regard to relative performance.

As illustrated in Figure \ref{fig:rl-analysis}, GRPO demonstrates a notable advantage over online RFT, underscoring the effectiveness of modifying positive and negative gradient coefficients. Moreover, GRPO+PS achieves better results than GRPO+OS, suggesting that employing fine-grained, step-sensitive gradient coefficients is beneficial. Additionally, our analysis extends to iterative RL, where we apply two rounds of iteration in our trials. The results depicted in Figure \ref{fig:iter-rl} reveal that iterative RL leads to pronounced improvements in performance, particularly evident in the initial iteration.

\subsubsection{Why RL Works?} This study applies reinforcement learning on a selected subset of instruction tuning data, resulting in marked performance gains over the instruction-tuned baseline. To further elucidate the reasons for RL's efficacy, we assess Pass@K and Maj@K accuracy metrics for both the Instruct and RL-optimized models across two benchmarks. As indicated in Figure \ref{fig:maj-pass}, RL mainly increases Maj@K, while leaving Pass@K relatively unchanged. This suggests that RL strengthens the overall output distribution robustness, meaning that the primary gain comes from increasing the frequency with which the correct answer appears among the TopK candidates, rather than by fundamentally enhancing baseline capabilities. Likewise, \citep{wang2023making} reported a \textbf{misalignment problem} in SFT models on reasoning tasks, and demonstrated that reasoning ability can be bolstered through targeted preference alignment strategies \citep{yuan2023rrhf,song2023preference,wang2023making}.

\subsubsection{How to Achieve More Effective RL?}
\label{sec:effec_rl}
We show that RL delivers substantial performance in mathematical reasoning benchmarks and introduce a unified perspective for characterizing various canonical training frameworks. Under this unified view, different approaches are interpreted as variants or reductions of RL-based methods. Equation \ref{eq:objective} outlines the three central ingredients: Data Source, Algorithm, and Reward Function. We propose several directions for further work related to each component.

\paragraph{Data Source} The data source forms the foundational input for any training method. In RL, this primarily refers to unlabeled questions paired with responses generated by the policy model. Our present work leverages only those questions drawn from instruction tuning and applies a straightforward nucleus sampling technique to sample responses. This constraint might explain why the observed RL gains are largely confined to Maj@K improvement. Moving forward, we aim to assess our RL method using question prompts that are out-of-distribution and to combine this with more sophisticated sampling (decoding) techniques, including tree-search based approaches \citep{yao2023tree}. Furthermore, advancements in \textbf{efficient inference techniques} \citep{xia-etal-2023-speculative,leviathan2023fast,kwon2023efficient,xia2024unlocking}, which critically influence how thoroughly policy models can be explored, will also be a key focus.

\paragraph{Algorithms} Algorithms leverage data and reward feedback to compute the gradient coefficients necessary for updating model parameters. Referring to Equation \ref{eq:objective}, prevailing approaches largely place complete \textbf{TRUST} in the reward function to adjust the conditional likelihood of generating specific tokens. Nonetheless, the dependability of the reward signal cannot be guaranteed across all scenarios, especially when dealing with highly intricate tasks. For instance, even the meticulously curated PRM800K datasets \citep{lightman2023let}, annotated by skilled professionals, still manifest about 20\% incorrect labels\footnote{\url{https://github.com/openai/prm800k/issues/12\#issuecomment-1728491852}}. Consequently, there is a need to investigate reinforcement learning methods that remain effective under noisy or imperfect reward signals. We anticipate that alignment strategies such as \textbf{WEAK-TO-STRONG} \citep{burns2023weak} will instigate a significant paradigm shift in learning algorithm design.

\paragraph{Reward Function} The reward function provides the essential feedback required during training. In the context of RL, this function typically takes the form of a neural reward model. We identify three pivotal research directions for reward models: 1) \textbf{Improving the generalization capacity of reward models.} For RL to foster substantive advancements in LLM performance rather than merely consolidating existing distributions, the reward model must be adept at handling both out-of-distribution inputs and sophisticated decoding scenarios; 2) \textbf{Quantifying and representing the uncertainty of the reward model.} Capturing uncertainty may serve as an interface between underperforming reward models and robust weak-to-strong alignment techniques; 3) \textbf{Developing efficient, high-quality process reward models} that deliver detailed and process-oriented feedback signals, thus supporting complex reasoning tasks \citep{lightman2023let,wang2023math}.

\section{Conclusion, Limitation, and Future Work}

This work introduces DeepSeekMath, which establishes new performance standards among open-source models on the MATH competition benchmark, narrowing the gap with proprietary models. DeepSeekMath~builds upon DeepSeek-Coder-v1.5 7B and is further trained with 500B tokens, including a substantial 120B tokens of mathematics-related data extracted from Common Crawl. Our comprehensive ablation analyses reveal that web pages provide considerable value as sources for mathematical content, whereas data from arXiv do not contribute as effectively as anticipated. We propose a new reinforcement learning strategy, Group Relative Policy Optimization (GRPO), a variation on Proximal Policy Optimization (PPO), which not only enhances mathematical reasoning but also reduces memory requirements. Experimental evidence demonstrates that GRPO delivers tangible benefits, even when DeepSeekMath-Instruct 7B already achieves top-tier benchmark results. Additionally, we establish a consolidated framework for situating these techniques and outline several avenues for advancing reinforcement learning in this context.

Despite these achievements, DeepSeekMath~still exhibits comparatively limited abilities in tackling geometry problems and theorem-proving tasks relative to proprietary counterparts. Specifically, dry-run tests indicate the model struggles with triangle and ellipse questions, suggesting possible biases or inadequacies in training and fine-tuning data curation. Furthermore, constrained by its architecture, DeepSeekMath~does not match GPT-4’s proficiency in few-shot learning scenarios—whereas GPT-4 benefits from few-shot prompts, our model’s performance is essentially unchanged between zero-shot and few-shot conditions. Moving forward, our plan is to refine the data selection process to compile a richer, higher-quality pre-training corpus. Moreover, we intend to pursue the promising research directions discussed in Section \ref{sec:effec_rl} to further advance reinforcement learning methods for large language models.

\end{document}